\section{기호와 단위 빠른 참조}
\label{sec:notation-checklist}

여기서부터는 본문을 이어 읽는 장이 아니라, 헷갈리는 기호를 다시 확인하기 위한 부록이다.
본문을 처음부터 끝까지 한 번에 이해하려고 하면 기호가 헷갈릴 수 있다.
특히 전기와 기계, 로봇 제어에서 비슷한 글자를 서로 다른 뜻으로 쓰는 경우가 많다.
이 부록은 새로운 이론을 추가하기 위한 장이 아니라, 본문을 읽다가 기호와 단위가 흔들릴 때 돌아와 확인하기 위한 빠른 참조표이다.
특히 $q$, $C$, $V$, $x$, $e$, $f$처럼 여러 분야에서 반복해서 등장하는 글자는 반드시 문맥과 아래첨자를 함께 확인해야 한다.
이 부록에서는 $q(t)$와 $Q(s)$를 전하로, $\bm{q}$와 $q_i$를 관절좌표로 구분한다.
또한 $\mathrm{C}$는 쿨롱 단위이고 $C$는 커패시턴스이며, $C_e$는 목표점 기준 오차 순응성이다.
$V(s)$는 전압이지만, $V_x(s)$와 $V_e(s)$처럼 아래첨자가 붙으면 속도의 라플라스 변환으로 읽는다.
$E(t)$는 에너지로, $E(s)$는 오차 $e(t)$의 라플라스 변환으로 쓰일 수 있으므로 문맥을 함께 확인한다.
회로 장에서 아래첨자 없이 쓰는 $X$는 리액턴스를 뜻할 수 있고, $X(s)$처럼 괄호와 라플라스 변수 $s$가 붙으면 위치 $x(t)$의 라플라스 변환으로 읽는다.

\subsection{초심자용 용어 길찾기}
\label{sec:beginner-glossary-roadmap}

다음 표는 새로운 이론을 더하는 표가 아니라, 본문을 읽다가 용어가 흔들릴 때 어디로 돌아가면 좋은지 알려 주는 색인이다.
처음 읽을 때는 모든 항목을 외우려고 하기보다, 헷갈리는 단어를 만났을 때 한 줄 의미를 보고 해당 절이나 표로 돌아가면 된다.

\begin{table}[H]
  \centering
  \caption{초심자용 로봇공학 용어 길찾기: 기구학과 궤적}
  \label{tab:beginner-glossary-kinematics}
  \small
  \setlength{\tabcolsep}{3pt}
  \renewcommand{\arraystretch}{1.14}
  \begin{tabularx}{\linewidth}{@{}>{\raggedright\arraybackslash}p{0.22\linewidth}>{\raggedright\arraybackslash}X>{\raggedright\arraybackslash}p{0.31\linewidth}@{}}
    \toprule
    용어 & 처음에는 이렇게 읽기 & 돌아갈 곳과 주의점 \\
    \midrule
    링크, 관절, 자유도 & 링크는 단단한 몸체, 관절은 허용된 움직임, 자유도는 독립적으로 정할 수 있는 움직임 개수이다 & 절~\ref{sec:robotics-links-joints-variables}, 표~\ref{tab:joint-variable-derivatives}. 자유도가 많으면 선택지도 많지만 해석도 쉬워진다는 뜻은 아니다 \\
    관절 변수와 미분 & $\bm{q}$는 관절 위치, $\dot{\bm{q}}$, $\ddot{\bm{q}}$, $\dddot{\bm{q}}$는 관절 속도, 가속도, 저크이다 & 표~\ref{tab:joint-variable-derivatives}. 전하 $q(t)$와 관절좌표 $\bm{q}$를 문맥으로 구분한다 \\
    관절공간, 작업공간, 작업영역 & 관절공간은 모터 각도들의 세계이고, 작업공간은 손끝 위치처럼 사람이 보고 싶은 세계이다 & 절~\ref{sec:robotics-joint-task-workspace}, 표~\ref{tab:joint-task-space}. 작업영역은 그중 실제로 도달 가능한 집합이다 \\
    정기구학(FK) & 관절값에서 손끝 위치나 포즈를 계산하는 일이다 & 절~\ref{sec:robotics-forward-kinematics}, 식~\eqref{eq:fk-minimal}--\eqref{eq:two-link-fk-y}. 관절값 하나를 정하면 이상적인 강체 모델의 손끝 위치가 정해진다 \\
    역기구학(IK) & 원하는 손끝 목표에서 필요한 관절값을 찾는 일이다 & 절~\ref{sec:robotics-inverse-kinematics}. 해가 없거나, 하나이거나, 여러 개일 수 있으며 여러 해는 자세 branch와 effort 차이로 이어진다 \\
    자코비안 & 관절을 조금 움직이면 손끝이 어느 방향으로 얼마나 움직이는지 알려 주는 순간 번역표이다 & 절~\ref{sec:robotics-jacobian-small-motion}, 식~\eqref{eq:jacobian-small-motion}, \eqref{eq:two-link-jacobian}. 큰 이동 전체를 푸는 공식이 아니라 현재 자세 근처의 선형 근사이다 \\
    속도기구학 & 작은 변화 식을 시간으로 나누어 관절속도와 손끝속도를 연결하는 관점이다 & 절~\ref{sec:robotics-velocity-kinematics}, 식~\eqref{eq:jacobian-velocity-foundation}. 가속도에는 $\dot{\bm J}\dot{\bm q}$ 항이 추가될 수 있다 \\
    특이점, 조건수, 조작성, DLS & 어떤 자세에서는 특정 손끝 방향이 만들기 어려워지고, 역문제가 작은 오차에도 예민해진다 & 절~\ref{sec:robotics-singularity-conditioning}, 식~\eqref{eq:two-link-manipulability}, \eqref{eq:dls-velocity-ik}. DLS 감쇠는 물리 감쇠비가 아니라 역기구학 계산 완화이다 \\
    경로, 궤적, 저크 프로파일 & 경로는 어디를 지나는가이고, 궤적은 언제 얼마나 빠르게 지나는가까지 포함한다 & 절~\ref{sec:robotics-path-trajectory}, \ref{sec:robotics-trajectory-profiles}, 표~\ref{tab:path-trajectory-profile}. 사다리꼴과 S-curve는 시간 스케일링을 읽는 말이다 \\
    \bottomrule
  \end{tabularx}
\end{table}

\begin{table}[H]
  \centering
  \caption{초심자용 로봇공학 용어 길찾기: 임피던스와 접촉}
  \label{tab:beginner-glossary-control-contact}
  \small
  \setlength{\tabcolsep}{3pt}
  \renewcommand{\arraystretch}{1.14}
  \begin{tabularx}{\linewidth}{@{}>{\raggedright\arraybackslash}p{0.22\linewidth}>{\raggedright\arraybackslash}X>{\raggedright\arraybackslash}p{0.31\linewidth}@{}}
    \toprule
    용어 & 처음에는 이렇게 읽기 & 돌아갈 곳과 주의점 \\
    \midrule
    자코비안 전치와 힘--관절 effort 변환 & 손끝 힘을 각 관절이 나눠 맡을 일반화 effort 쪽으로 옮길 때 $\bm{J}^{\mathsf T}$가 등장한다. 회전 관절에서는 이 값이 토크이다 & 절~\ref{sec:robotics-force-torque-jacobian-transpose}, 식~\eqref{eq:jacobian-transpose-foundation}, 표~\ref{tab:jacobian-recap-for-impedance}. 좌표계와 힘 부호를 함께 확인한다 \\
    평형점, 목표점, anchor point & 가상 스프링이 묶인 기준 위치이다. 힘이 균형을 이루는 곳으로 읽으면 된다 & 식~\eqref{eq:steady_position_under_force}, \eqref{eq:virtual_spring_damper}. 목표점 오프셋과 실제 벽 침투 깊이를 섞지 않는다 \\
    강성 & 위치가 목표에서 벗어날수록 되돌리려는 힘이 얼마나 빨리 커지는지를 나타낸다 & 식~\eqref{eq:virtual_spring_damper}. 접촉 방향 강성을 키우면 정적 힘이 커질 수 있지만 토크 포화와 안전 한계도 같이 본다 \\
    감쇠와 감쇠비 & 감쇠는 속도에 저항해 흔들림을 줄이는 효과이고, 감쇠비는 질량과 강성까지 함께 본 흔들림 모양의 기준이다 & 식~\eqref{eq:damping_from_ratio}. DLS damping, virtual wall damping, 물리 감쇠비를 같은 값으로 읽지 않는다 \\
    임피던스 제어와 어드미턴스 제어 & 임피던스는 운동에서 힘 반응을 설계하는 관점이고, 어드미턴스는 힘을 받아 목표 운동을 바꾸는 관점이다 & 식~\eqref{eq:task_impedance}, 표~\ref{tab:impedance-implementation-ladder}. 실제 로봇 인터페이스가 힘/토크 명령인지 위치 명령인지 먼저 본다 \\
    관절 임피던스와 직교 임피던스 & 관절 기준으로 가상 스프링을 거는지, 손끝 작업공간 기준으로 거는지의 차이이다 & 식~\eqref{eq:joint_impedance}, \eqref{eq:cartesian_to_joint_torque}. 직교 임피던스에서는 자코비안 조건과 방향별 단위를 조심한다 \\
    가상 벽, 침투 깊이, 접근 속도 & 정해 둔 경계를 넘으면 계산식으로 밀어내는 교육용 벽 모델이다 & 표~\ref{tab:contact-symbols}, 식~\eqref{eq:virtual-wall-penetration}--\eqref{eq:virtual-wall-force}. 가상 벽 힘은 실측 접촉력이 아니라 정의한 식으로 계산한 신호이다 \\
    position actuator와 effort proxy & 위치 액추에이터는 목표 관절 위치를 받는 인터페이스이고, effort proxy는 구동기 부담 비교용 대리 지표이다 & 표~\ref{tab:artifact-theory-bridge}, \ref{tab:lab-observation-guide}. 실측 모터 전류나 힘 센서값으로 읽지 않는다 \\
    6차원 포즈, 트위스트, 렌치 & 자세까지 포함하면 위치만이 아니라 회전 오차, 6차원 속도, 힘과 모멘트 묶음이 필요하다 & 절~\ref{sec:notation-6d-pose-impedance}, 표~\ref{tab:pose-impedance-notation}. 이 논문의 Lab04는 완전한 6차원 포즈 임피던스 검증이 아니다 \\
    \bottomrule
  \end{tabularx}
\end{table}

\subsection{영역 표기 먼저 구분하기}

\begin{table}[H]
  \centering
  \caption{시간영역, 라플라스 영역, 주파수 응답의 빠른 구분}
  \small
  \setlength{\tabcolsep}{4pt}
  \begin{tabularx}{\linewidth}{@{}>{\raggedright\arraybackslash}p{0.21\linewidth}>{\raggedright\arraybackslash}X>{\raggedright\arraybackslash}X@{}}
    \toprule
    표기 & 쉬운 의미 & 읽을 때의 질문 \\
    \midrule
    $x(t)$, $v(t)$ & 시간에 따라 변하는 실제 신호 & 지금 순간의 움직임이나 전압을 보는가? \\
    $X(s)$, $V(s)$ & 라플라스 영역의 신호 & 과도응답과 안정성까지 함께 보는가? \\
    $Z(\jj\omega)$, $H(\jj\omega)$ & 정현파 주파수 응답 & 특정 주파수에서 크기와 위상을 보는가? \\
    $s=\sigma+\jj\omega$ & 감쇠와 진동을 함께 담은 복소 변수 & $s$를 시간 단위 초와 혼동하지 않았는가? \\
    \bottomrule
  \end{tabularx}
\end{table}

\subsection{전기 시스템 기호}

\begin{table}[H]
  \centering
  \caption{전기 시스템에서 자주 쓰는 기호}
  \small
  \setlength{\tabcolsep}{4pt}
  \begin{tabularx}{\linewidth}{@{}>{\raggedright\arraybackslash}p{0.15\linewidth}>{\raggedright\arraybackslash}p{0.27\linewidth}>{\raggedright\arraybackslash}p{0.14\linewidth}>{\raggedright\arraybackslash}X@{}}
    \toprule
    기호 & 의미 & 대표 단위 & 주의할 점 \\
    \midrule
    $v(t)$, $V(s)$ & 전압 & $\mathrm{V}$ & 전기적 effort이다 \\
    $i(t)$, $I(s)$ & 전류 & $\mathrm{A}$ & 전기적 flow이다 \\
    $q(t)$, $Q(s)$ & 전하 & $\mathrm{C}$ & $q_i$, $\bm{q}$는 관절각이므로 문맥으로 구분한다 \\
    $R$ & 저항 & $\Omega$ & 에너지를 열로 소산한다 \\
    $L$ & 인덕턴스 & $\mathrm{H}$ & 전류 변화에 저항하고 자기장에 에너지를 저장한다 \\
    $C$ & 커패시턴스 & $\mathrm{F}$ & 전압/전하 상태를 전기장 에너지로 저장하고 전압의 급격한 변화에 저항한다. 전하 단위 $\mathrm{C}$와 구분한다 \\
    $Z_e$ & 전기 임피던스 & $\Omega$ & $V/I$이며, 페이저나 라플라스 문맥에서 읽는다 \\
    \bottomrule
  \end{tabularx}
\end{table}

\subsection{기계 시스템 기호}

\begin{table}[H]
  \centering
  \caption{기계 시스템에서 자주 쓰는 기호}
  \small
  \setlength{\tabcolsep}{4pt}
  \begin{tabularx}{\linewidth}{@{}>{\raggedright\arraybackslash}p{0.16\linewidth}>{\raggedright\arraybackslash}p{0.26\linewidth}>{\raggedright\arraybackslash}p{0.15\linewidth}>{\raggedright\arraybackslash}X@{}}
    \toprule
    기호 & 의미 & 대표 단위 & 주의할 점 \\
    \midrule
    $x(t)$ & 위치 & $\mathrm{m}$ & 스프링 변위나 로봇 끝단 위치로 쓰인다 \\
    $\dot{x}(t)$, $V_x(s)$, $V_e(s)$ & 속도 & $\mathrm{m/s}$ & 아래첨자가 붙은 $V_x$, $V_e$는 전압이 아니라 속도이다 \\
    $e$, $\bm{e}$ & 위치 오차 & $\mathrm{m}$ & 보통 $e=x-x_d$, $\bm{e}=\bm{x}-\bm{x}_d$이다. $e^{st}$의 $e$는 자연상수이다 \\
    $E(s)$ & 오차의 라플라스 변환 & $\mathrm{m}$ & 제어공학의 신호 단위 관례로 적는다. 엄밀한 적분 변환 단위는 별도 초 단위를 포함할 수 있다 \\
    $f(t)$, $F(s)$ & 힘 & $\mathrm{N}$ & 기계적 effort이다 \\
    $m$ & 질량 & $\mathrm{kg}$ & 운동 에너지를 저장하고 속도 변화에 저항한다 \\
    $b$ & 점성 감쇠계수 & $\mathrm{N\,s/m}$ & 속도에 비례해 에너지를 소산한다 \\
    $k$ & 스프링 강성 & $\mathrm{N/m}$ & 탄성 퍼텐셜 에너지를 저장하고 위치가 평형점에서 벗어나는 것을 되돌린다 \\
    $Z_m$ & 기계 임피던스 & $\mathrm{N\,s/m}$ & 힘/속도 관계이며 힘/위치가 아니다 \\
    $C_e$ & 오차 순응성 & $\mathrm{m/N}$ & 힘이 목표점 기준 위치 오차를 얼마나 만드는지 나타낸다 \\
    $Y_e$ & 오차 어드미턴스 & $\mathrm{m/(N\,s)}$ & 오차 속도/힘 관계이고, 고정 목표점에서는 기계 어드미턴스 $Y_m=V_x/F_{\mathrm{ext}}$처럼 읽을 수 있다 \\
    $Y_m$ & 기계 어드미턴스 & $\mathrm{m/(N\,s)}$ & 절대 속도/힘 관계이다. 목표점이 움직이면 $Y_e=V_e/F_{\mathrm{ext}}$와 구분한다 \\
    \bottomrule
  \end{tabularx}
\end{table}

\subsection{2차 시스템과 로봇 제어 기호}

\begin{table}[H]
  \centering
  \caption{2차 시스템과 로봇 제어에서 자주 쓰는 기호}
  \small
  \setlength{\tabcolsep}{4pt}
  \begin{tabularx}{\linewidth}{@{}>{\raggedright\arraybackslash}p{0.16\linewidth}>{\raggedright\arraybackslash}p{0.26\linewidth}>{\raggedright\arraybackslash}p{0.15\linewidth}>{\raggedright\arraybackslash}X@{}}
    \toprule
    기호 & 의미 & 대표 단위 & 주의할 점 \\
    \midrule
    $\omega$ & 각주파수 & $\mathrm{rad/s}$ & Hz 변환은 $f=\omega/(2\pi)$이다 \\
    $\omega_n$ & 고유진동수 & $\mathrm{rad/s}$ & 시스템이 자연스럽게 흔들리려는 빠르기이다 \\
    $\zeta$ & 감쇠비 & 없음 & 감쇠계수 하나가 아니라 $m,b,k$의 조합이다 \\
    $m_d$, $\bm{M}_d$ & 희망 관성, 목표 모델의 관성 & $\mathrm{kg}$ & 제어기가 만들고 싶은 가상 관성이다. 실제 링크 질량이나 끝단 유효질량과 구분한다 \\
    끝단 유효질량, 겉보기 관성 & 현재 자세와 힘 방향에서 손끝이 무겁게 또는 가볍게 느껴지는 값 & $\mathrm{kg}$ & 로봇 구조와 자세가 정하는 값이며, 사용자가 고르는 $m_d$나 $m_{\mathrm{nom}}$과 다르다 \\
    $m_{\mathrm{nom}}$ & 교육용 명목 질량 & $\mathrm{kg}$ & $d_d$ 계산을 일관되게 하기 위한 기준값이다. 실제 끝단 유효질량이나 안전성 보증으로 읽지 않는다 \\
    $d_d$, $\bm{D}_d$ & 희망 감쇠, 감쇠계수 & $\mathrm{N\,s/m}$ & 1축 또는 대각 근사에서는 $d_d=2\zeta\sqrt{m_dk_d}$로 계산한다. 교육용 명목 질량을 쓰면 $d_d=2\zeta\sqrt{m_{\mathrm{nom}}k_d}$이며, 실제 유효질량과 다르면 감쇠비는 보증값이 아니라 튜닝 기준이다 \\
    $k_d$, $\bm{K}_d$ & 희망 강성 & $\mathrm{N/m}$ & 병진 기준이다. 실제 끝단 느낌은 자코비안과 액추에이터 인터페이스의 영향을 받는다 \\
    $\bm{K}\succeq0$, $\bm{K}\succ0$ & 양의 준정부호, 양의 정부호 & 없음 & $\succeq0$은 에너지가 음수가 되지 않는다는 뜻이고, $\succ0$은 모든 방향에서 닫힌 그릇처럼 복원력이 있다는 뜻이다 \\
    $\Delta x$ & 허용 변위 & $\mathrm{m}$ & 원하는 접촉력에서 강성 출발값을 잡을 때 쓰는 설계 허용량이다. $k_d\approx F/\Delta x$로 읽으며, 실제 침투 $\delta$나 목표 오프셋 $\delta_d$와 구분한다 \\
    $\delta_d$ & 목표점 오프셋 & $\mathrm{m}$ & 벽 뒤에 일부러 둔 목표 평형점의 거리이다. 접촉력을 만들기 위한 설계량이며 실제 침투 깊이가 아니다 \\
    $\delta$ & 벽 침투 깊이 & $\mathrm{m}$ & 현재 끝단이 벽을 넘은 양이다. 가상 벽에서는 $\delta=\max(0,x-x_{\mathrm{wall}})$로 두며 $\delta_d$와 별도이다 \\
    $\bm{x}$ & 작업공간 위치 & $\mathrm{m}$ & 이 글에서는 주로 병진 위치 $\R^3$에 집중한다 \\
    $\bm{q}$, $q_i$ & 관절공간 좌표 & 회전 관절은 $\mathrm{rad}$, 직선 관절은 $\mathrm{m}$ & 전기 전하 $q(t)$와 문맥으로 구분한다. 하나의 벡터 안에서도 관절 종류에 따라 성분 단위가 달라질 수 있다 \\
    $\bm{J}_v$ & 병진 자코비안 & 회전 관절 열은 $\mathrm{m/rad}$, 직선 관절 열은 $\mathrm{m/m}$ & 병진 위치 기준이다. $\dot{\bm{x}}=\bm{J}_v(\bm{q})\dot{\bm{q}}$로 작은 운동을 연결한다 \\
    조작성(manipulability) & 손끝을 여러 방향으로 움직이기 쉬운 정도 & 없음 & 작아지면 특이점에 가까워져 역기구학이 예민해질 수 있다 \\
    나쁜 조건수 상태(ill-conditioned) & 계산이 오차에 예민한 상태 & 없음 & 작은 위치 오차나 센서 노이즈가 큰 관절 변화로 증폭될 수 있다 \\
    $\bm{\tau}$ & 일반화 관절 effort & 회전 관절은 $\mathrm{N\,m}$, 직선 관절은 $\mathrm{N}$ & 병진 힘은 $\bm{\tau}=\bm{J}_v^{\mathsf{T}}\bm{f}$로 관절 effort와 연결된다. 회전 관절 성분은 토크, 직선 관절 성분은 축 방향 힘이다 \\
    \bottomrule
  \end{tabularx}
\end{table}

관절공간 임피던스를 직접 쓸 때는 단위가 달라진다.
관절 강성은 보통 $\mathrm{N\,m/rad}$, 관절 감쇠는 $\mathrm{N\,m\,s/rad}$로 읽는 편이 안전하다.
따라서 $\bm{K}_d$, $\bm{D}_d$가 작업공간 값인지 관절공간 값인지 항상 먼저 확인해야 한다.

\subsection{병진 위치와 6차원 포즈 표기 구분}
\label{sec:notation-6d-pose-impedance}

본문의 대부분은 손끝 병진 위치 $\bm{x}=[x\ y\ z]^{\mathsf T}\in\R^3$를 기준으로 설명했다.
즉, 로봇 손끝이 공간에서 어디에 있는지를 보는 설명이다.
하지만 실제 로봇 작업에서는 손끝이 어디에 있는지만큼, 공구가 어느 방향을 보고 있는지도 중요할 수 있다.
예를 들어 드라이버를 나사에 맞추거나, 문손잡이를 돌리거나, 컵을 기울이지 않고 옮기는 작업에서는 자세(orientation)를 빼면 문제가 완전히 설명되지 않는다.

이때 위치와 자세를 합쳐 포즈(pose)라고 부른다.
자세는 회전행렬 $R\in SO(3)$로, 위치와 자세를 함께 담은 포즈는 보통 $T\in SE(3)$로 쓴다.
처음 읽을 때는 $R$은 손끝 좌표축이 어느 방향을 보는지, $T$는 위치와 방향을 함께 담은 좌표계 표지판이라고 생각하면 된다.
롤--피치--요 세 각도로 자세를 표현할 수도 있지만, 특정 자세 근처에서 특이점이 생길 수 있으므로 정밀한 제어식에서는 회전행렬, 사원수, 또는 회전 오차 벡터를 따로 쓴다.

초심자용으로는 다음 구분만 먼저 붙잡으면 된다.
병진 위치는 \(\bm{x}\in\mathbb{R}^3\)이지만, 6차원 포즈는 보통
\[
  T
  =
  \begin{bmatrix}
    R & \bm{p} \\
    \bm{0}^{\mathsf T} & 1
  \end{bmatrix}
\]
처럼 위치 \(\bm{p}\)와 자세 \(R\)를 함께 담는 표현이다.
반면 트위스트는 포즈 자체가 아니라 순간 속도
\[
  \bm{V}
  =
  \begin{bmatrix}
    \bm{v} \\
    \bm{\omega}
  \end{bmatrix}
\]
이고, 렌치는 힘과 모멘트
\[
  \bm{w}
  =
  \begin{bmatrix}
    \bm{f} \\
    \bm{\mu}
  \end{bmatrix}
\]
이다.
즉 포즈는 현재의 위치와 방향을 적은 사진이고, 트위스트는 그 사진이 지금 어떻게 움직이는지를 적은 속도이다.

\begin{table}[H]
  \centering
  \caption{병진 위치 임피던스와 6차원 포즈 임피던스 표기 구분}
  \label{tab:pose-impedance-notation}
  \small
  \setlength{\tabcolsep}{4pt}
  \renewcommand{\arraystretch}{1.12}
  \begin{tabularx}{\linewidth}{@{}>{\raggedright\arraybackslash}p{0.20\linewidth}>{\raggedright\arraybackslash}p{0.31\linewidth}>{\raggedright\arraybackslash}X@{}}
    \toprule
    표기 & 쉬운 의미 & 읽을 때의 주의점 \\
    \midrule
    $\bm{x}\in\R^3$ & 손끝 병진 위치 & 이 논문의 주된 작업공간 변수이다. 회전 자세는 포함하지 않는다 \\
    $R\in SO(3)$ & 손끝 자세, 즉 공구 좌표축의 방향 & 롤--피치--요 숫자 3개와 같다고 단순화하면 특이점과 좌표계 문제가 생길 수 있다 \\
    $T\in SE(3)$ & 위치와 자세를 합친 포즈 & 손끝이 어디에 있고 어느 방향을 보는지를 함께 담는다 \\
    롤--피치--요와 각속도 & 자세를 세 각도로 적는 편리한 좌표와 실제 순간 회전 속도는 구분한다 & 롤--피치--요의 시간미분은 특이점 근처에서 불안정할 수 있고, 일반적으로 $\bm{\omega}$와 같은 벡터가 아니다. 6차원 트위스트에는 Euler angle rate가 아니라 같은 좌표계에서 표현한 각속도 $\bm{\omega}$를 넣는다 \\
    $\bm{\omega}$, $\dot{\bm{\omega}}$ & 각속도, 각가속도 & 병진 속도 $\dot{\bm{x}}$와 단위가 다르다. 각각 $\mathrm{rad/s}$, $\mathrm{rad/s^2}$로 읽는다 \\
    $\bm{V}$ & 트위스트, 병진 속도와 각속도를 묶은 6차원 속도 & 선속도와 각속도를 같은 벡터에 넣지만 단위가 섞여 있다 \\
    $\bm{w}$ & 렌치, 힘과 모멘트를 묶은 6차원 힘 & 병진 힘 $\bm{f}$와 회전 모멘트 $\bm{\mu}$를 함께 다룬다 \\
    $\bm{J}_g$ & 기하학적 자코비안 & $\bm{V}=\bm{J}_g\dot{\bm{q}}$처럼 관절속도를 6차원 손끝 속도로 옮긴다 \\
    $\bm{\tau}=\bm{J}_g^{\mathsf T}\bm{w}$ & 렌치를 관절 토크로 옮기는 관계 & $\bm{V}$, $\bm{w}$, $\bm{J}_g$가 같은 기준 좌표계와 기준점에서 표현되어야 한다 \\
    \bottomrule
  \end{tabularx}
\end{table}

따라서 6차원 포즈 임피던스는 단순히 $\bm{x}$를 6개 숫자로 늘리는 일이 아니다.
병진 오차 $\bm{x}-\bm{x}_d$에 더해 회전 오차를 어떤 방식으로 정의할지 정해야 하고, 힘뿐 아니라 모멘트까지 포함한 렌치를 써야 한다.
회전 오차는 단순히 $R-R_d$처럼 행렬을 빼서 끝나는 값으로 생각하기보다, 현재 자세에서 목표 자세까지 어느 축으로 얼마나 돌려야 하는지를 나타내는 작은 회전 벡터로 먼저 읽는 편이 안전하다.
한 축만 도는 아주 쉬운 경우에는 yaw rate가 \(\omega_z\)처럼 보일 수 있다.
하지만 일반적인 3차원 자세에서는 롤--피치--요 각도의 시간미분이 각속도 \(\bm{\omega}\) 자체와 같지 않고, 선택한 자세 표현에 따른 변환을 거쳐야 한다.
각가속도도 Euler 각도의 두 번째 미분을 그대로 넣는 값으로 단순화하면 안 되고, 선택한 자세 표현과 좌표계 변환을 확인해야 한다.
또 병진 강성의 단위는 $\mathrm{N/m}$이지만 회전 강성은 보통 $\mathrm{N\,m/rad}$로 읽는다.
둘을 한 행렬에 넣을 수는 있지만, 단위와 좌표계가 다르다는 사실을 잊으면 안 된다.
특히 $\bm{V}$, $\bm{w}$, $\bm{J}_g$는 같은 좌표계와 같은 기준점에서 표현되어야 한다.
공구 좌표계(tool frame)에서 측정한 렌치를 기준 좌표계(base frame) 자코비안에 바로 넣으면 방향과 모멘트 기준점이 틀어질 수 있다.

이 논문에서 Lab04를 병진 위치 데모로 제한한 이유도 여기에 있다.
현재 Lab04 설명은 손끝의 $x,y,z$ 위치, 계산된 가상 벽 힘, DLS 목표 오프셋, position actuator effort 지표를 읽는 훈련이다.
표~\ref{tab:pose-impedance-notation}의 6차원 표기는 나중에 자세와 모멘트까지 다루는 로봇 제어를 읽기 위한 안내판이지, 현재 Lab04가 완전한 6차원 포즈 임피던스를 검증한다는 뜻이 아니다.

\subsection{힘 부호 빠른 구분}

\begin{table}[H]
  \centering
  \caption{임피던스 제어에서 힘 기호를 읽는 법}
  \small
  \setlength{\tabcolsep}{4pt}
  \begin{tabularx}{\linewidth}{@{}>{\raggedright\arraybackslash}p{0.22\linewidth}>{\raggedright\arraybackslash}X>{\raggedright\arraybackslash}p{0.28\linewidth}@{}}
    \toprule
    기호 & 의미 & 부호/크기 주의 \\
    \midrule
    $f_{\mathrm{ext}}$ & 환경이 로봇에 가하는 외력 & 부호 있는 힘 \\
    $f_{\mathrm{cmd}}$ & 로봇 제어기가 만들려고 하는 명령힘 & 부호 있는 힘 \\
    $f_{\mathrm{wall}}$ & 벽이 로봇에 가하는 힘 & 부호 있는 힘 \\
    $|f_{\mathrm{contact}}|$ & 접촉력의 크기 & 방향보다 크기를 볼 때 사용 \\
    \bottomrule
  \end{tabularx}
\end{table}

\subsection{구현 신호 빠른 구분}

\begin{table}[H]
  \centering
  \caption{임피던스 관계를 실제 로봇 명령으로 옮길 때의 구분}
  \small
  \setlength{\tabcolsep}{4pt}
  \begin{tabularx}{\linewidth}{@{}>{\raggedright\arraybackslash}p{0.25\linewidth}>{\raggedright\arraybackslash}p{0.34\linewidth}>{\raggedright\arraybackslash}X@{}}
    \toprule
    로봇 인터페이스 & 제어기가 만드는 신호 & 읽을 때의 기준 \\
    \midrule
    힘/토크 명령 가능 & 위치/속도 오차에서 $\bm{f}_{\mathrm{cmd}}$ 또는 $\bm{\tau}$를 계산 & 임피던스식 구현으로 읽는다. 단, 토크 한계와 동역학 보상이 실제 응답을 제한한다 \\
    위치 명령만 가능 & 측정 또는 추정 외력에서 목표 위치/속도 변화를 계산 & 어드미턴스식 외부 루프 또는 목표 위치 보정으로 읽는다. Lab04 가상 벽은 실측 외력 대신 계산된 가상 벽 신호를 쓰는 교육용 목표 후퇴 근사이며, 완전한 운영공간 임피던스 보증은 아니다 \\
    \bottomrule
  \end{tabularx}
\end{table}

\subsection{읽을 때마다 확인할 질문}

수식을 만났을 때 다음 질문을 차례대로 던지면 훨씬 덜 헷갈린다.

\begin{enumerate}
  \item 지금 식은 시간영역 식인가, 라플라스 영역 식인가, 정현파 정상상태의 페이저 식인가?
  \item 입력과 출력은 무엇인가? 힘에서 위치로 보는가, 힘에서 속도로 보는가, 운동에서 힘으로 보는가?
  \item 정상상태를 말하고 있는가, 과도응답을 말하고 있는가?
  \item 힘은 부호 있는 힘인가, 접촉력의 크기인가?
  \item 위치 좌표인가, 속도 좌표인가, 관절공간인가, 작업공간인가?
  \item 강성, 감쇠, 관성 중 어떤 요소가 에너지를 저장하고 어떤 요소가 에너지를 소산하는가?
  \item 같은 기호가 다른 장에서 다른 뜻으로 쓰이지 않는가? 특히 $q$, $C$, $V$는 문맥을 확인한다.
  \item 이 설명은 실제 로봇 성능 보증인가, 교육용 모델에서 개념을 확인하기 위한 단순화인가?
\end{enumerate}

수식을 만났을 때 이 질문들은 잠깐 방향을 잡아 주는 작은 지도에 가깝다.
지금 식이 시간응답을 말하는지, 주파수응답을 말하는지, 입력과 출력이 무엇인지 알면 계산보다 먼저 그림이 잡힌다.
처음에는 목록을 하나씩 확인해도 좋고, 익숙해지면 필요한 질문만 빠르게 꺼내 쓰면 된다.
헷갈리는 순간에 이 중 하나만 붙잡아도 다음 문장이 훨씬 편하게 이어진다.
